\chapter{พื้นฐานโรคใบพืชทุเรียนและวิทยาศาสตร์เฉดสีของรอยโรค}

\section{ความสำคัญของการตรวจวินิจฉัยโรคใบพืชทุเรียน}
ทุเรียน (\textit{Durio zibethinus} L.) จัดเป็นราชาแห่งผลไม้และเป็นพืชเศรษฐกิจเชิงยุทธศาสตร์ที่สร้างรายได้มหาศาลให้แก่ประเทศไทย ประสิทธิภาพการสังเคราะห์แสง การสะสมสารอาหารคาร์โบไฮเดรต และการเจริญเติบโตเพื่อสร้างผลผลิตคุณภาพสูง ล้วนขึ้นอยู่กับความสมบูรณ์และพื้นที่ใบสีเขียวของเรือนพุ่มเป็นหลัก อย่างไรก็ตาม สภาพภูมิอากาศแบบร้อนชื้นในแถบภาคตะวันออกและภาคใต้ เป็นสภาวะแวดล้อมที่เอื้อต่อการแพร่ระบาดของเชื้อรา เชื้อสาเหตุโรคพืช และศัตรูพืชปากดูดนานาชนิด ซึ่งการระบาดที่รวดเร็วอาจทำลายพื้นที่ใบของต้นทุเรียนจนส่งผลให้ต้นทรุดโทรม ใบร่วงหล่น และผลผลิตเสียหายอย่างสิ้นเชิง

การตรวจวินิจฉัยรอยโรคในระยะเริ่มต้นด้วยเทคโนโลยีคอมพิวเตอร์วิทัศน์ (Computer Vision) และปัญญาประดิษฐ์เชิงแสง จึงเป็นกลไกสำคัญในการตัดวงจรการระบาด ช่วยให้เกษตรกรสามารถบริหารจัดการสวนได้อย่างแม่นยำและทันท่วงที

\section{การจำแนกกลุ่มโรคและศัตรูพืชสำคัญ 6 กลุ่ม}
ในการพัฒนาระบบปัญญาประดิษฐ์ DurianLeaf AI คณะผู้วิจัยได้กำหนดการจำแนกกลุ่มโรคและศัตรูพืชหลักออกเป็น 6 กลุ่มอาการตามเกณฑ์กีฏวิทยาและโรคพืชวิทยา ดังแสดงในตารางที่ \ref{tab:durian_disease_colorimetry}

\begin{table}[h]
\centering
\caption{คุณลักษณะทางแสงและพิกัดสีมาตรฐานของรอยโรคใบพืชทุเรียน 6 กลุ่ม}
\label{tab:durian_disease_colorimetry}
\small
\begin{tabular}{|c|l|c|c|c|c|l|}
\hline
\textbf{ลำดับ} & \textbf{ชื่อโรคและศัตรูพืช} & \textbf{รหัสสี HEX} & $L^*$ & $a^*$ & $b^*$ & \textbf{ลักษณะอาการรอยโรค} \\
\hline
1 & ใบปกติสมบูรณ์ & \texttt{\#2E7D32} & 48.2 & -22.4 & 28.6 & เขียวมรกตเข้ม สม่ำเสมอ ไม่พบคราบหรือแผล \\
2 & ใบไหม้ไฟทอปธอร่า & \texttt{\#3E2723} & 26.5 & 8.4 & 12.1 & แผลฉ่ำน้ำสีน้ำตาลเข้มถึงดำ ขอบแผลไหม้ \\
3 & ราใบติด (Rhizoctonia) & \texttt{\#6D4C41} & 38.1 & 5.6 & 14.8 & น้ำตาลปนเทา คล้ายน้ำร้อนลวก มีใยราคลุม \\
4 & แอนแทรคโนส & \texttt{\#BF360C} & 44.8 & 16.2 & 29.5 & วงแหวนซ้อน concentric สีน้ำตาลส้ม \\
5 & ราสนิมและจุดสนิมสาหร่าย & \texttt{\#D84315} & 56.3 & 18.5 & 42.0 & ตุ่มนูนส้มอมแดง ผิวสัมผัสคล้ายกำมะหยี่ \\
6 & ไรแดงและเพลี้ยไก่แจ้ & \texttt{\#9E9D24} & 62.4 & -2.1 & 18.7 & จุดประสีบรอนซ์เงิน ใบซีดด่าง แห้งกร้าน \\
\hline
\end{tabular}
\end{table}

\subsection{โรคใบไหม้ไฟทอปธอร่า (Phytophthora Leaf Blight)}
เกิดจากเชื้อราน้ำโอโอไมซีต (\textit{Phytophthora palmivora}) อาการเริ่มแรกจะปรากฏเป็นจุดแผลฉ่ำน้ำสีน้ำตาลอ่อน ขอบแผลไม่เรียบ แผลจะขยายตัวลุกลามอย่างรวดเร็วเป็นสีน้ำตาลเข้มจนถึงดำ (\texttt{\#3E2723}) ในสภาพความชื้นสัมพัทธ์สูงกว่า $90\%$ เนื้อเยื่อบริเวณแผลจะเน่าเละ ขอบแผลมีลักษณะฉ่ำน้ำและแห้งกรอบเป็นวงกว้าง ส่งผลให้ใบหลุดร่วงอย่างรวดเร็ว

\subsection{โรคราใบติด (Rhizoctonia Web Blight)}
เกิดจากเชื้อรา \textit{Rhizoctonia solani} อาการบนแผ่นใบมีลักษณะเด่นคล้ายถูกน้ำร้อนลวก เนื้อเยื่อใบเปลี่ยนเป็นสีน้ำตาลปนเทา (\texttt{\#6D4C41}) เชื้อราสร้างเส้นใยขาวนวลหรือสีน้ำตาลถักทอยึดแผ่นใบที่อยู่ติดกันให้ติดกันเป็นแพ เมื่อแผลแห้งกรอบ ใบจะแห้งตายและร่วงค้างอยู่บนกิ่งไม่ยอมหล่นลงดิน

\subsection{โรคแอนแทรคโนส (Anthracnose)}
เกิดจากเชื้อรา \textit{Colletotrichum gloeosporioides} มักปรากฏแผลรูปทรงกลมหรือรูปรีขอบหยัก แผลมีสีน้ำตาลอ่อนปนส้ม (\texttt{\#BF360C}) โดยมีจุดเด่นคือรอยวงแหวนซ้อนกันเป็นชั้น (Concentric Rings หรือ Zonate Pattern) บริเวณขอบแผลจะมีสีเข้มกว่าเนื้อเยื่อแผลด้านใน กลางแผลแห้งบางและอาจฉีกขาดกลายเป็นรูพรุน

\subsection{โรคราสนิมและจุดสนิมสาหร่าย (Algal Rust Spot)}
เกิดจากสาหร่ายกาฝากสีเขียวชนิด \textit{Cephaleuros virescens} มักพบบนใบแก่ที่ได้รับแสงแดดจัด แผลมีลักษณะเป็นตุ่มกลมนูนขนาดเล็ก $2 - 6\text{ mm}$ ผิวหน้าแผลมีขนละเอียดคล้ายกำมะหยี่สีส้มอิฐหรือสีสนิม (\texttt{\#D84315}) ค่าสี $b^*$ จะมีค่าบวกสูง ($>40$) เมื่อแก่ตัวลงจุดสาหร่ายจะเปลี่ยนเป็นสีน้ำตาลอมเทา

\subsection{ความเสียหายจากไรแดงและเพลี้ยไก่แจ้ (Mite and Psyllid Damage)}
เกิดจากการเข้าทำลายของศัตรูพืชปากดูด เช่น ไรแดงแอฟริกัน (\textit{Oligonychus biharensis}) และเพลี้ยไก่แจ้ทุเรียน (\textit{Pseudophacopteron canarium}) แมลงจะดูดกินน้ำเลี้ยงบริเวณหน้าใบและยอดอ่อน ทำให้คลอโรฟิลล์ถูกทำลายเป็นจุดประพร่างสีขาวซีดหรือสีบรอนซ์เงิน (\texttt{\#9E9D24}) ค่าความสว่าง $L^*$ จะพุ่งสูงขึ้นสู่ช่วง $60 - 65$ และค่าความอิ่มตัวสีลดลง แผ่นใบจะโค้งงอ ผิวใบสากกร้านและร่วงก่อนกำหนด

\section{เกณฑ์การจัดระดับความรุนแรงของโรค (Disease Severity Index)}
ระบบได้กำหนดมาตรฐานการจัดเกรดความรุนแรงของโรคออกเป็น 5 ระดับ (Grade 0 ถึง Grade 4) เพื่อนำไปประเมินดัชนีความรุนแรง (Disease Severity Index หรือ DSI) สำหรับการวางแผนพ่นสารชีวภัณฑ์หรือสารป้องกันกำจัดโรคพืช ดังแสดงในตารางที่ \ref{tab:durian_dsi_levels}

\begin{table}[h]
\centering
\caption{เกณฑ์การจัดระดับความรุนแรงของโรคใบพืชทุเรียน (DSI)}
\label{tab:durian_dsi_levels}
\begin{tabular}{|c|c|l|l|}
\hline
\textbf{ระดับ (Grade)} & \textbf{สัดส่วนพื้นที่แผล (\%)} & \textbf{นิยามทางพยาธิสภาพ} & \textbf{มาตรการจัดการแปลง} \\
\hline
Grade 0 & $0\%$ & ใบสมบูรณ์ แข็งแรง ไม่พบรอยโรค & ตรวจติดตามรายสัปดาห์ \\
Grade 1 & $1 - 10\%$ & ระยะเริ่มต้น พบรอยแผลกระจัดกระจาย & ตัดแต่งใบที่เป็นโรคทิ้ง ฉีดพ่นเชื้อราไตรโคเดอร์มา \\
Grade 2 & $11 - 25\%$ & ระยะลุกลาม แผลเริ่มเชื่อมต่อกัน & ใช้สารชีวภัณฑ์หรือคอปเปอร์ไฮดรอกไซด์ \\
Grade 3 & $26 - 50\%$ & ระยะรุนแรง แผ่นใบถูกทำลายเกินหนึ่งในสี่ & ใช้สารเคมีกำจัดเชื้อราเฉพาะกลุ่ม สลับกลไกออกฤทธิ์ \\
Grade 4 & $>50\%$ & ระยะวิกฤต ใบไหม้เกือบทั้งใบ แห้งร่วง & ตัดทำลายกิ่ง เก็บใบเผานอกแปลง หยุดการแพร่ระบาด \\
\hline
\end{tabular}
\end{table}

\begin{theorybox}
\textbf{การประเมินดัชนีความรุนแรงของโรค (Disease Severity Index: DSI)}\par
การคำนวณค่าดัชนีความรุนแรงของโรคในแปลงทุเรียนใช้สูตรมาตรฐานทางระบาดวิทยาพืช
\begin{equation}
  \text{DSI} = \frac{\sum_{i=0}^{4} (n_i \times i)}{N \times 4} \times 100\%
\end{equation}
เมื่อ $n_i$ คือจำนวนใบที่ถูกจัดอยู่ในระดับความรุนแรง $i$, $N$ คือจำนวนใบตัวอย่างทั้งหมดที่สุ่มตรวจ และ $4$ คือระดับความรุนแรงสูงสุด ค่า DSI จะถูกแสดงผลสดบนหน้าจอเพื่อช่วยให้เกษตรกรตัดสินใจระดับการใช้สารควบคุมโรคพืชได้อย่างคุ้มค่าและปลอดภัย
\end{theorybox}

\begin{promptbox}
\textbf{พรอมพ์วิศวกรรมสำหรับการสกัดรอยโรคและการวิเคราะห์พยาธิวิทยา (Prompt Eng-01)}\par
\texttt{"Act as a Senior Plant Pathologist and Computer Vision Specialist. Analyze the provided multi-channel durian leaf image within the designated Region of Interest (ROI). Segment the lesion pixels using CIE L*a*b* chromaticity bounds and texture contrast. Classify the pathogen into one of the 6 core classes: Healthy, Phytophthora Blight, Rhizoctonia Web Blight, Anthracnose, Algal Rust Spot, or Mite/Psyllid Damage. Calculate the precise Lesion Area Ratio, map to Disease Severity Index (DSI) Grade 0 to 4, and output the result along with immediate biological and chemical treatment recommendations."}
\end{promptbox}
